% !TITLE 山坡羊·潼关怀古
% !AUTHOR 张养浩
% !DYNASTY 元
% !GENRE 散曲
% !ORDER 2

\begin{poem}
峰峦如聚\textsuperscript{1}，波涛如怒\textsuperscript{2}，山河表里\textsuperscript{3}潼关\textsuperscript{4}路。
望西都\textsuperscript{5}，意踌躇\textsuperscript{6}。
伤心秦汉经行处\textsuperscript{7}，宫阙万间都做了土\textsuperscript{8}。
兴，百姓苦；亡，百姓苦。
\end{poem}

\begin{annotations}
\notetext{1}{峰峦如聚：山峰从四面八方聚集过来。形容潼关地势险要，群山环绕。}
\notetext{2}{波涛如怒：黄河的波涛像在发怒。潼关北临黄河，河水奔腾咆哮。}
\notetext{3}{山河表里：外面是黄河（表），里面是华山（里）。形容潼关依山傍河的地理形势。}
\notetext{4}{潼关：古代著名关隘，位于今陕西潼关县，是进入关中平原的咽喉要道。}
\notetext{5}{西都：指长安（今西安）。长安在潼关以西。}
\notetext{6}{踌躇：犹豫、徘徊。也指心潮起伏、百感交集。}
\notetext{7}{秦汉经行处：秦朝和汉朝曾经经营过的地方。长安是秦汉两代的都城。}
\notetext{8}{宫阙万间都做了土：千万间宫殿楼阁都化为了泥土。指朝代更替，昔日繁华荡然无存。}
\end{annotations}

\begin{appreciation}
《山坡羊·潼关怀古》是元代散曲里最有分量的一首之一。山坡羊是曲牌名，"潼关怀古"才是题目。作者张养浩和流落市井的曲家不同，官至礼部尚书。

"峰峦如聚，波涛如怒"——开头八个字就把潼关写活了：山峰从四面聚拢，黄河的浪像在发怒。一个"聚"字一个"怒"字，全是拟人：山不是立着，是聚过来；河不是流着，是吼着冲过来。"山河表里潼关路"——外黄河，内华山，一条险路锁住进关中的咽喉。

"望西都，意踌躇"——潼关西望就是长安，秦汉两朝的都城。"伤心秦汉经行处，宫阙万间都做了土"——他走的这条路秦汉帝王都走过，他们盖的万间宫阙如今都成了土。写到这里还只是寻常怀古：繁华终将成土，谁都会说。

但最后两句一出，前面全部作废："兴，百姓苦；亡，百姓苦。"朝代兴盛，百姓苦——大兴土木、服徭役；朝代灭亡，百姓苦——战乱、离散、饿殍。兴和亡本是反义词，落到百姓身上居然一个样。八个字，把三千年历史的账本一页翻到了底。

这不是书斋里的感慨。张养浩因为直言敢谏辞官归隐。天历二年关中大旱，朝廷起用他赈灾，他散尽家财就出发，路过潼关写下这首曲子，几个月后死在任上。他是走进过"百姓苦"里的人，所以这两句从他嘴里说出来，每个字都沉。屈原在《离骚》里说"哀民生之多艰"，一千七百年后，张养浩在潼关把它凝成了八个字。

我们书里的画面都浮起来了：《十五从军征》里亲人全没了的老兵；《兵车行》里"信知生男恶"的哭喊；《卖炭翁》里一车炭被抢走的老人。汉代、唐代、元代，换了一朝又一朝，路边的百姓没换过。

历史书里写的是英雄和王朝，兴亡的代价却从来是普通人扛的。哥哥不指望你背朝代年号，只希望你将来无论做什么工作，见到辛苦劳动的人，心里知道是他们扛着这个时代。别看不起任何一个用力活着的人。
\end{appreciation}
